\section*{Chapter \ref{ch:g4:clouds-rain-snow} --- Clouds, Rain and Snow}

\begin{solution}{exo:g4:clouds-rain-snow:1}
A \omterm{def:g4:clouds-rain-snow:cloud}{cloud} is made of a great number of tiny droplets of water, or tiny pieces
of ice, floating in the air. \omterm{def:g4:clouds-rain-snow:cloud}{Fog} is a \omterm{def:g4:clouds-rain-snow:cloud}{cloud} that touches the ground.
\end{solution}

\begin{solution}{exo:g4:clouds-rain-snow:2}
(a) \omterm{def:g4:clouds-rain-snow:cloud-families}{Stratus}. (b) \omterm{def:g4:clouds-rain-snow:cloud-families}{Cirrus}. (c) \omterm{def:g4:clouds-rain-snow:cloud-families}{Cumulus}. (d) \omterm{def:g4:clouds-rain-snow:cloud-families}{Cumulonimbus}.
\end{solution}

\begin{solution}{exo:g4:clouds-rain-snow:3}
(a) \omterm{def:g4:clouds-rain-snow:precipitation}{Hail}. (b) \omterm{def:g4:clouds-rain-snow:precipitation}{Rain}. (c) \omterm{def:g4:clouds-rain-snow:precipitation}{Snow}.
\end{solution}

\begin{solution}{exo:g4:clouds-rain-snow:4}
(a) A \omterm{def:g4:clouds-rain-snow:rain-gauge}{rain gauge}. (b) A thermometer. (c) A \omterm{def:g4:clouds-rain-snow:wind-vane}{wind vane}.
\end{solution}

\begin{solution}{exo:g4:clouds-rain-snow:5}
17 millimetres of \omterm{def:g4:clouds-rain-snow:precipitation}{rain} have fallen. $20 - 17 = 3$: 3 more millimetres
would bring the water to the mark 20.
\end{solution}

\begin{solution}{exo:g4:clouds-rain-snow:6}
The wettest day was Wednesday (12 millimetres). It did not \omterm{def:g4:clouds-rain-snow:precipitation}{rain} on Monday
and on Friday. Thursday had 7 millimetres.
\end{solution}

\begin{solution}{exo:g4:clouds-rain-snow:7}
The snowflakes form high in the cold \omterm{def:g4:clouds-rain-snow:cloud}{cloud}. Up in the \omterm{def:g1:landscapes-around-us:landform}{mountain} village the
air is freezing all the way down, so they arrive as \omterm{def:g4:clouds-rain-snow:precipitation}{snow}. Lower down, in
the \omterm{def:g1:landscapes-around-us:landform}{valley}, the air is warmer than \SI{0}{\celsius}: on the way down
through it the flakes melt, and they reach the ground as \omterm{def:g4:clouds-rain-snow:precipitation}{rain}.
\end{solution}

\begin{solution}{exo:g4:clouds-rain-snow:8}
The bar chart has seven bars of heights 3, 0, 1, 9, 6, 0 and 1 (no bar on
Tuesday and Saturday). In the week, $3 + 0 + 1 + 9 + 6 + 0 + 1 = 20$
millimetres fell. Thursday brought 9 of the 20 millimetres: $\frac{9}{20}$
of the week's \omterm{def:g4:clouds-rain-snow:precipitation}{rain}.
\end{solution}

\begin{solution}{exo:g4:clouds-rain-snow:9}
Walls and trees shelter the ground near them from the \omterm{def:g4:clouds-rain-snow:precipitation}{rain}, especially when
the wind blows the \omterm{def:g4:clouds-rain-snow:precipitation}{rain} sideways, and drops can drip from leaves into the
gauge: near them the gauge would catch too little, or too much. The
thermometer has to measure the temperature of the air; in full sun it would be
warmed by the sunshine itself and would read too high. The white slatted box
keeps the sun off and lets the air through.
\end{solution}

\begin{solution}{exo:g4:clouds-rain-snow:10}
A \omterm{def:g4:clouds-rain-snow:wind-vane}{wind vane} points to where the wind comes from: the wind comes from the
west. It is going towards the east.
\end{solution}

\begin{solution}{exo:g4:clouds-rain-snow:11}
One millimetre of \omterm{def:g4:clouds-rain-snow:precipitation}{rain} on one square metre is one litre. On one square
metre, 8 millimetres give 8 litres. On 500 square metres:
$500 \times 8 = 4000$ litres. With watering cans of 10 litres:
$4000 \div 10 = 400$ watering cans.
\end{solution}
